\section{Scope \& Methodology}\label{sec:method}
\subsection{Subject \& Scope}
\begin{center}\begin{tabularx}{\linewidth}{@{}l X@{}}
\toprule\rowcolor{TableHead}\textcolor{white}{\bfseries Item} & \textcolor{white}{\bfseries Detail} \\
\midrule
System          & \reportsystem \\
\rowcolor{TableAlt}
Investigation type & \fillme{e.g. Read-only filesystem/artifact review -- discovery only, nothing executed or modified} \\
Data sources    & \fillme{e.g. Deployed application source tree; filesystem mtimes; no live network capture} \\
\bottomrule\end{tabularx}\end{center}

\subsection{Approach}
\begin{enumerate}
  \item \textbf{Artifact discovery.} \fillme{How suspicious files/backdoors were located -- e.g.
        pattern search for exec/eval sinks, unusual file extensions or permissions, filename
        anomalies.}
  \item \textbf{Characterisation.} \fillme{How each artifact was analysed to determine what it
        does, without executing it.}
  \item \textbf{Timeline reconstruction.} \fillme{How dates/ordering were established -- e.g.
        filesystem mtimes -- and the confidence caveats that come with that method.}
  \item \textbf{Impact \& IOC synthesis.} Findings were compiled into a consolidated inventory,
        a timeline, and a set of indicators of compromise for detection/hunting.
\end{enumerate}

\begin{notebox}
\small\textbf{Caveat.} \fillme{Note any limitation, e.g. "filesystem mtimes are
attacker-influenceable and bound presence, not necessarily the moment of first abuse."}
\end{notebox}
